% Independent supplementary calculation.  No main file is edited here.
\subsection*{The coefficient-two, order-three Kummer component}

\paragraph{KM.1 (Proved): a full profinite definition with no logarithm.}
Let $F=\widehat F(x,y)$, $N_n=\ker(F\to C_n)$, where $x$ maps
to a chosen generator and $y$ to one.  A free profinite basis is
\[
 x^n,\quad y_j=x^jyx^{-j}\quad(0\leq j<n).
\]
For $f\in F'$, write its image in $N_n^{\rm ab}$ as
\[
 0\cdot[x^n]+\sum_{j=0}^{n-1}m_j[y_j].
\]
The coefficient of $x^n$ is zero and $\sum_jm_j=0$, because
$f$ has zero exponent sums in $F^{\rm ab}$.
Define $\rho_n(f)=m_0$.
For normalized automorphisms $\Phi_g(x)=x^{\lambda_g}$,
$\Phi_g(y)=f_g^{-1}y^{\lambda_g}f_g$ with $f_g\in F'$, this is
a continuous crossed homomorphism:
\[
 \rho_n(gh)=\rho_n(g)+\lambda_g\rho_n(h).
\]
Indeed, $N_n$ is preserved, and its abelianization sends
$[y_j]$ to $\lambda_g[y_{\lambda_gj}]$ and $[x^n]$ to
$\lambda_g[x^n]$.  Conjugation by a power of $x^n$ does not
affect this statement.  Since multiplication of normalized pairs
gives $f_{gh}=f_g\Phi_g(f_h)$, the displayed formula follows.
Thus the restriction to the multiplier-one group is an additive
character with values in $\widehat{\mathbb Z}$.

This is the $\tau_n^{(0)}$ definition in Ihara's primary
RIMS-1224 manuscript, pp.3 and 6 (the preprint of
\emph{On beta and gamma functions associated with the
Grothendieck--Teichm\"uller group. II}, J. reine angew. Math.
527 (2000), 1--11). Theorem 1, manuscript p.4, writes
$(n^{-r})^{\sigma-1}=e[r\tau_n^{(0)}(\sigma)]$.
Thus, in that paper's monodromy convention, it is the negative
of the exponent $k_n$ defined by
$\sigma(n^{1/m})/n^{1/m}=\zeta_m^{k_n(\sigma)}$ for the positive roots.
Replacing both free generators by their inverses also changes the
Schreier coefficient $m_0$ to $-m_0$.
The even half modulo two used below is invariant under either change;
no identification of the integral signs in two different markings
is being made.

\paragraph{KM.2 (Proved): an order-48 formula for $\rho_3$ modulo four.}
Let
\[
 R=\mathbb Z_2[\omega],\qquad \omega^2+\omega+1=0,
 \qquad (t,u)(t',u')=(t+ut',uu')
\]
in $R\rtimes\langle\omega\rangle$.  Evaluate at
$X=(0,\omega)$ and $Y=(1,1)$, and write
$f(X,Y)=(t_f,1)$ for $f\in F'$.  Then
\[
 \rho_{3,2}(f)=\frac{\operatorname{Tr}_{R/\mathbb Z_2}(t_f)}3.
 \tag{KM.2}
\]
In fact $y_j$ maps to $(\omega^j,1)$, so
$t_f=m_0+m_1\omega+m_2\omega^2$.
Using $m_0+m_1+m_2=0$ and
$\operatorname{Tr}(1)=2$, $\operatorname{Tr}(\omega)
=\operatorname{Tr}(\omega^2)=-1$ gives
$\operatorname{Tr}(t_f)=3m_0$.
Modulo four the target is the explicitly finite group
\[
 ((\mathbb Z/4)[\omega],+)\rtimes C_3
\]
of order $48$.  If $t_f=a+b\omega$, the formula reads
$\rho_{3,2}\equiv3(2a-b)\pmod4$.

Equivalently, with $[x,y]=xyx^{-1}y^{-1}$ and
$f\equiv A_f(x,y)[x,y]\pmod{F''}$, one has
\[
 t_f=(\omega-1)A_f(\omega,1),\qquad
 \rho_{3,2}(f)=\frac13\operatorname{Tr}
       ((\omega-1)A_f(\omega,1)).
\]
This is the coefficient-two component of the displayed formula
in Schneps (2006), p.657.  The Schreier proof fixes the additive
normalization independently of the evident typographical errors
in the target and root-of-unity subscript printed beside that formula.
The finite target retains order three; it is not a pro-two
quotient of $F$.

\paragraph{KM.3 (Proved in the standard gentle framework): the half-character is integral.}
Suppose $\Phi=(1,f)$ is an actual gentle element, with the standard
two-cycle and three-cycle equations, so that
\[
 \Phi(z)=f(z,x)z f(z,x)^{-1},\qquad z=(xy)^{-1}.
 \tag{KM.3a}
\]
Then
\[
 t_f\in\mathbb Z_2\subset R,\qquad
 \rho_{3,2}(f)\in2\mathbb Z_2,\qquad
 \frac{\rho_{3,2}(f)}2=\frac{t_f}3.
 \tag{KM.3b}
\]
In particular its half modulo two is computed in the order-twelve
group $\mathbb F_4\rtimes C_3\simeq A_4$: the translation of
$f(X,Y)$ is either $0$ or $1$, and is exactly this bit.
Since two is a unit in every odd-adic factor, the full
$\rho_3$ consequently lies in $2\widehat{\mathbb Z}$.
Its unique half is a continuous additive character on the
multiplier-one gentle group, by KM.1 and injectivity of
multiplication by two on $\widehat{\mathbb Z}$.

For completeness, (KM.3a) follows directly from the two three-cycle
equations available for an actual gentle element.
Put $g=f(z,x)$ and $h=f(y,z)$. The multiplier-one equation gives
$ghf=1$. Right composition with $c=(-1,1)$ gives the normalized
pair $(-1,f)$, whose three-cycle, with $m=-1$, is
$gz^{-1}hy^{-1}fx^{-1}=1$. Substituting $h=g^{-1}f^{-1}$ yields
$gz^{-1}g^{-1}=xf^{-1}yf=\Phi(z^{-1})$, proving (KM.3a).
This uses the stated group structure of the actual gentle group,
not the pentagon or a lifting assertion.

Here are the remaining details, to separate this elementary evenness from
any extra gamma multiplication relation. First evaluate at
$X_0=(0,\omega)$, $Y_0=(1,\omega^{-1})$.  The product
$X_0Y_0$ is a translation.  The conjugator in (KM.3a) is also a
translation, since $f\in F'$, and therefore acts trivially on
$X_0Y_0$.  If $f(X_0,Y_0)=(v,1)$, the equality
$\Phi(xy)=xy$ in this target gives
$(1-\omega)v=0$.  Since $1-\omega$ is a unit in $R$, $v=0$.
Since $[X_0,Y_0]=(\omega-1,1)$, this proves
\[
 A_f(\omega,\omega^{-1})=0.
 \tag{KM.3c}
\]
The two-cycle implies $A_f(U,V)=A_f(V,U)$ in the free
metabelian module.  Now return to $X=(0,\omega)$,
$Y=(1,1)$ and $Z=(XY)^{-1}=(-1,\omega^{-1})$.
The three-cycle is
$f(Z,X)f(Y,Z)f(X,Y)=1$.  Its first factor is one by
(KM.3c); its other translations are, respectively,
\[
 (1-\omega^{-1})A_f(1,\omega^{-1})=-\overline{t_f},
 \qquad t_f.
\]
Consequently $t_f=\overline{t_f}$, where the bar exchanges
$\omega$ and $\omega^{-1}$.  Its fixed ring is $\mathbb Z_2$.
Formula (KM.2) now proves (KM.3b).  This proof uses no pentagon
and does not identify either modular frame sign with this character.

\paragraph{KM.4 (Proved): the exact remaining type of comparison.}
For multiplier one the intrinsic metabelian action on the
commutator module is multiplication by
\[
 B_f(U,V)=1+(U-1)(V-1)A_f(U,V).
\]
To check the sign, use affine coordinates with linear parts $U,V$:
$\Phi(y)=f^{-1}yf$ changes its translation by $(V-1)f$,
and the translation of $[x,\Phi(y)]$ changes by
$(U-1)(V-1)A_f$ times that of $[x,y]$.
The same plus convention is used in the proof of NTY (2010),
Proposition 4.4, p.443.  At the diagonal order-three component,
\[
 B_f(\omega,\omega)
   =1+ (\omega-1)^2A_f(\omega,\omega)
   =1-3\omega A_f(\omega,\omega).
\]
The five-point audit's affine comparison of this value with the
Gaussian frame concerns this \emph{diagonal} component, whereas
(KM.2) concerns $(\omega,1)$.  Equating the resulting bit with
$\rho_{3,2}/2$ is a further assertion, not a consequence of the
definition of $\rho_3$ or of its evenness.

Ihara II, introduction pp.1--2 and Section 2, makes the relevant
source boundary explicit: full profinite pentagon gives
$B(s,t)=\Gamma(s)\Gamma(t)/\Gamma(s+t)$, but Anderson's
multiplication formula is imposed as an additional condition
defining $\mathrm{GT}_A$.  Its logarithmic-derivative version
likewise defines $\mathrm{GT}_{A'}$.  Proposition 2 proves
$\mathrm{GT}_K\subseteq\mathrm{GT}_{A'}$, and Proposition 3
identifies $(A')$ with $(D\log A)$ and yields
$\mathrm{GT}_A\subseteq\mathrm{GT}_{A'}$.
The final Remark explicitly warns that the derivation has a
nonconstant kernel from mixed prime components.  Thus these
opened results do not license replacing a coefficient-two,
order-three unit or its sign by logarithmic Kummer data.
\paragraph{KM.5 (Proved, coefficient-ring statement only): factorization and logarithmic data do not determine the sign.}
This is a countermodel to a proposed deduction between coefficient-ring
identities, not a constructed element of GT.
Let $\Lambda=\mathbb Z_2[[\widehat{\mathbb Z}]]$ and let
$\Lambda_3=\mathbb Z_2[[\mathbb Z_3]]$ be its subalgebra from
the pro-three direct factor of the exponent group.  The invariant
derivation $D(x^a)=a_2x^a$, where $a_2$ is the two-adic component
of $a$, vanishes identically on $\Lambda_3$.

There is an integral idempotent $e\in\Lambda_3$ supported on
characters of exact order three.  Explicitly, in the quotient
group ring $\mathbb Z_2[C_{3^n}]$ with generator $v$, for $n\geq1$
its image is
\[
 e_n=\frac1{3^{n-1}}\sum_{j=0}^{3^{n-1}-1}v^{3j}
       -\frac1{3^n}\sum_{j=0}^{3^n-1}v^j.
\]
These are compatible idempotents: the first summand projects
onto the characters of order dividing three and the second onto
the trivial character.  Put $\Gamma=1-2e$.  Then
\[
 \Gamma^2=1,\quad D\Gamma=0,\quad \Gamma(1)=1,
 \quad \Gamma(\omega)=\Gamma(\omega^{-1})=-1.
\]
Moreover $\Gamma(U^{-1})=\Gamma(U)^{-1}$, since inversion fixes
$e$.  Define in the completed two-variable coefficient ring
\[
 B(U,V)=\Gamma(U)\Gamma(V)\Gamma(UV)^{-1}.
\]
It is a normalized symmetric multiplicative two-cocycle and
satisfies $B(U,U^{-1})=1$.  Nevertheless
\[
 B(\omega,\omega)=-1,
 \qquad D_VB(\omega,1)=0.
\]
The first assertion follows from $\omega^2=\omega^{-1}$;
the second follows from $D\Gamma=0$.  Thus the diagonal unit
sign cannot be recovered from the off-diagonal logarithmic
derivative using only these identities.  To apply a stronger
identity to the actual GT image, that additional identity must
be proved or cited in its full profinite, mixed-component domain.

The usual integrality constraint on the metabelian coefficient
does not remove this particular coefficient-ring example.
The odd-character decomposition of the completed group algebra
is integral over $\mathbb Z_2$.  On any component where the
pro-three character of $U$ or of $V$ is trivial, $B(U,V)=1$.
On every remaining component, both $U-1$ and $V-1$ are units:
their nontrivial roots of odd order remain nontrivial in residue
characteristic two.  Thus an integral $A$ exists with
$B=1+(U-1)(V-1)A$, taking $A=0$ on the former components.
It has $A(\omega,1)=0$ and
$A(\omega,\omega)=2/(3\omega)$.
No profinite free-group lift satisfying the pentagon is claimed.

